\section{Dataset-as-Input: matriz de datos, mask y contrato de predicción}

Para que todo el conjunto de datos entre al modelo, primero hay que reescribir la notación de aprendizaje supervisado (supervised-learning notation) como una matriz unificada. NPT trata las filas (rows) como datapoints y las columnas (columns) como atributos (attributes) con semántica compartida; la etiqueta (label) también ocupa una columna de atributo. La mask matrix determina qué entradas (entries) se observan y cuáles son incógnitas que el modelo debe predecir.

\lecturefigure{slide-13-entire-dataset-input.jpg}{Primer elemento: tomar el entire dataset, es decir, todos los datapoints, como entrada del modelo.}{Video oficial de Stanford Online 00:24:00--00:24:13.}

Sea
\[
X\in\mathbb{R}^{n\times d},\qquad M\in\{0,1\}^{n\times d}.
\]

\begin{itemize}
  \item $n$: número de datapoints en el forward actual; con datos pequeños puede ser el conjunto de datos completo, con datos grandes suele ser un minilote (mini-batch);
  \item $d$: número de atributos, que incluye las input features y puede incluir también las columnas objetivo (target columns);
  \item $X_{ij}$: el $j$-ésimo atributo del $i$-ésimo datapoint;
  \item $M_{ij}=1$: la entrada está enmascarada (mask) y debe predecirse; $M_{ij}=0$: la entrada es visible.
\end{itemize}

\lecturefigure{slide-14-datapoint-attention.jpg}{Segundo elemento: aplicar self-attention entre datapoints.}{Video oficial de Stanford Online 00:24:14--00:24:27.}

En un lote en tiempo de prueba (test-time batch), la attention puede conectar a la vez filas train--train, test--test y train--test. Aquí «train/test» es un papel determinado por la visibilidad del objetivo (target visibility): las filas de entrenamiento pueden exponer parte de las etiquetas, mientras que las etiquetas de las filas de prueba deben ocultarse por completo. El modelo ve la misma matriz, pero sabe por la mask qué entradas pueden servir como contexto.

\lecturefigure{slide-15-masking-objective-overview.jpg}{Tercer elemento: usar un masking-based objective para aprender de forma conjunta la reconstrucción del target y de las features.}{Video oficial de Stanford Online 00:24:28--00:24:41.}

Se definen las entradas masked y observed:
\[
X^{\mathcal{M}}=\{X_{ij}\mid M_{ij}=1\},\qquad
X^{\mathcal{O}}=\{X_{ij}\mid M_{ij}=0\}.
\]
NPT aprende
\[
p\!\left(X^{\mathcal{M}}\mid X^{\mathcal{O}},M\right),
\]
y produce $\widehat X\in\mathbb{R}^{n\times d}$. Esta representación convierte la clasificación, la regresión, la imputation, la self-supervision y la semi-supervision en «cambiar la posición de la mask».

\begin{importantbox}{La mask es la interfaz de la tarea}
Una misma architecture puede expresar problemas distintos mediante $M$: ocultar la columna de etiquetas de las filas de prueba produce supervised prediction; ocultar al azar entradas de features produce imputation/self-supervision; ocultar varias columnas objetivo produce multi-target prediction. Que la interfaz del modelo esté unificada no significa que la distribución de los datos y la loss de las distintas tareas queden unificadas automáticamente.
\end{importantbox}

\subsection{Notación completa: por qué también hay que dar la mask como entrada al modelo}

Las tres figuras anteriores presentaron el dataset, la datapoint attention y el masking objective; esta subsección los reúne en una interfaz de entrada sin ambigüedades. Al leer la figura de notación completa, conviene distinguir primero «cuál es el valor» de «si el valor es visible»: un mismo cero puede ser una observación válida o un marcador de posición para un dato faltante, y con solo mirar $X$ no es posible distinguirlos. Incrustar $M$ junto con $X$ equivale a indicarle explícitamente al modelo el epistemic status de cada valor: si es evidence conocida o una query por predecir; esto también determina en qué entradas debe aplicarse la loss.

\lecturefigure{slide-16-full-dataset-mask-notation.jpg}{Contrato de entrada completo: dataset $X$, mask $M$ y objetivo $p(X^{\mathcal{M}}\mid X^{\mathcal{O}})$.}{Video oficial de Stanford Online 00:25:06--00:27:15.}

\begin{knowledgebox}{Lectura de la figura: los tres colores y los signos de interrogación}
Cada fila es un datapoint y cada columna es un atributo de semántica fija. El signo de interrogación no significa «valor igual a cero», sino una entrada desconocida con $M_{ij}=1$. El modelo debe aprovechar a la vez las demás features de la misma fila, las features de otras filas y los targets de otras filas cuya exposición está permitida para reconstruir las posiciones marcadas con signo de interrogación.
\end{knowledgebox}

\begin{warningbox}{El límite infranqueable del test-label leakage}
El stochastic target masking puede conservar parte de las etiquetas en las filas de entrenamiento, lo que ayuda al modelo a aprender el lookup entre filas; las etiquetas reales de las filas de prueba nunca pueden usarse como entrada. Si al construir el lote o la mask se expone una test label, un resultado vistoso solo demuestra que el pipeline filtra la respuesta.
\end{warningbox}

\teachervoice{El ponente enuncia primero los tres elementos en conjunto y luego los desarrolla uno por uno: el conjunto de datos como entrada, la datapoint attention y el masking objective. Este orden nos recuerda que no conviene fijarse solo en el architecture block; lo que realmente define a NPT también incluye la input semantics y la training task.}

\subsection{Resumen de la sección}

NPT usa $(X,M)$ para representar de manera unificada la evidence y la query. Las filas son datapoints, las columnas son atributos de semántica compartida y la mask determina la tarea de predicción. Este contrato permite que los puntos train/test entren juntos a la attention y, al mismo tiempo, exige proteger estrictamente los test targets.
